\section{Introduction}
\label{sec:introduction}

The primitive residual program has reached two different analytic
boundaries.  The first is a conductor boundary: after exact
shared-factor rectangularization, two complementary estimates for the
same conductor cell can produce a fixed power saving.  The second is a
Poisson boundary: the four centered divisibility channels have periods
\([u,v]\), \(u\), \(v\), and \(1\), so a single channel is completed by
smoothness only when its period is shorter than the orbit length.

TPC-20 derived the exact four-channel formula and completed both single
channels under
\begin{equation}
 H\max(u,v)\le JX^{-\eta_0}.
 \label{eq:intro-both-short}
\end{equation}
TPC-22 closed the resulting joint short--short packet by exact
conductors, and TPC-23 moved the common cutoff from
\(R=X^{1/2-\delta+o(1)}\) to \(S>R\), still subject to
\eqref{eq:intro-both-short}
\citep{WangTPC20,WangTPC22,WangTPC23}.  The longer single channel, not
the conductor tensor, stops that centered-block argument at \(S<J\).

This paper moves only one divisor leg.  Let
\begin{equation}
 R\le S<T\le U,\qquad
 S=X^{s+o(1)},\quad T=X^{t+o(1)},\quad 0<s\le t,
 \label{eq:intro-ST}
\end{equation}
where \(U=X^{O(1)}\) is the finite target horizon.  Put
\[
 a(v)=-\mu(v)\log v,
 \qquad
 b_R(u)=a(u)-\lambda'_R(u),
\]
with \(\lambda'_R(u)=0\) for \(u>R\).  The new divisor region is the
incremental rectangle
\begin{equation}
 u\le S,\qquad S<v\le T.
 \label{eq:intro-mixed-shell}
\end{equation}
It is essential that the second leg is a complete M\"obius shell,
not an isolated long modulus or an arbitrary dyadic subset.

\subsection{The calibrated increment is the stable object}

Let \(P_{m,W}(j)\) be the locally calibrated residual prefix defined in
\cref{eq:calibrated-prefix}.  Its exact finite identity is
\begin{equation}
 P_{m,W}(j)
 =\sum_{\substack{w\le W\\w\mid mj+h}}b_R(w)-\delta_H(m),
 \qquad
 \delta_H(\ell d_0)\ll_{\epsilon,h}\frac{X^\epsilon}{L}.
 \label{eq:intro-calibrated-prefix}
\end{equation}
Because the finite model ends at \(R\), the increment is
\begin{equation}
 L_{m;S,T}(j):=P_{m,T}(j)-P_{m,S}(j)
 =\sum_{\substack{S<w\le T\\w\mid mj+h}}a(w).
 \label{eq:intro-exact-increment}
\end{equation}
We study the separated packet
\begin{align}
 \cI_{S,T}:={}&
 \sum_{m,n}x_my_n\mathfrak m(m,n)
 \sum_{\substack{j\in\Z\\(j,H)=1}}F(j/J)
 P_{m,S}(j)L_{n;S,T}(j).
 \label{eq:intro-calibrated-packet}
\end{align}
Its transpose is identical after exchanging the row labels.

The pointwise product in \eqref{eq:intro-calibrated-packet} has the
exact expansion
\begin{align}
 P_{m,S}(j)L_{n;S,T}(j)
 ={}&
 \sum_{\substack{u\le S\\u\mid mj+h}}b_R(u)
 \sum_{\substack{S<v\le T\\v\mid nj+h}}a(v)
 -\delta_H(m)L_{n;S,T}(j).
 \label{eq:intro-joint-minus-drift}
\end{align}
This identity is the first structural result.  If one starts instead
from two bare centered divisor sums, a long \(v\)-single Poisson
channel survives and is multiplied by a merely logarithmically
calibrated short density.  In the genuine prefix increment, the scalar
calibration terms restore the omitted channels and cancel that long
single exactly.  What remains is one joint congruence lattice and the
already power-saving row drift in \eqref{eq:intro-calibrated-prefix}.

\subsection{The asymmetric conductor rectangle}

For one compatible pair, write
\[
 g=(u,v),\qquad u=gc,\qquad v=gd,\qquad q=[u,v]=g\,c\,d.
\]
The TPC-22 character coordinates and weighted exact-conductor large
sieve extend to the rectangle \(u\le S\), \(v\le T\).  If
\(Q=X^{\beta+o(1)}\), \(J=X^{1-\beta+o(1)}\), then every joint
nonzero conductor cell has one of two complementary bounds.  Their
three-cell ledger gives
\begin{equation}
 \boxed{
 \eta_{\rm asym}(\beta;s,t)=
 \min\left\{
 1-\frac{3\beta}{2},
 \frac{\beta+1-2s-2t}{2},
 \frac{3-\beta-s-5t}{4}
 \right\}.}
 \label{eq:intro-asym-saving}
\end{equation}
Thus the joint nonzero discrepancy has a fixed power saving whenever
\(\eta_{\rm asym}>0\) with a fixed margin.  No hypothesis \(T<J\) or
\(L>2T\) occurs in this conductor estimate.

\subsection{The physical zero mode and the mass ledger}

The induced-Gauss and Ramanujan-frequency estimates explicitly exclude
additive frequency zero.  For an off-diagonal row pair put
\(\Delta=m-n\).  The zero-frequency kernel of the joint term in
\eqref{eq:intro-joint-minus-drift} is
\begin{equation}
 K^0_{S,T}(m,n)=
 \sum_{\substack{u\le S,\ S<v\le T\\
                  (u,mH)=(v,nH)=1}}
 \frac{b_R(u)a(v)}{uv}
 (u,v)\one_{(u,v)\mid\Delta}.
 \label{eq:intro-zero-kernel}
\end{equation}
Choose \(0<\kappa<s\).  Write \(u=gc\), \(v=gd\), with
\(g\mid\Delta\).  If \(g\le SX^{-\kappa}\), then
\(S/g\ge X^\kappa\).  For fixed \(g,c\), the inner \(d\)-sum is a
difference of uniformly calibrated M\"obius sums at \(T/g\) and
\(S/g\); their common main constant cancels.  If
\(g>SX^{-\kappa}\), then the off-diagonal mask gives
\begin{equation}
 \sum_{\substack{g\mid\Delta\\g>SX^{-\kappa}}}\frac1g
 \le \frac{X^\kappa}{S}\tau(|\Delta|).
 \label{eq:intro-large-g}
\end{equation}
Consequently, for every fixed \(A>0\),
\begin{equation}
 \boxed{
 K^0_{S,T}(m,n)
 \ll_{A,\kappa}
 (\log X)^{-A}+X^{-s+\kappa+o(1)}.}
 \label{eq:intro-zero-bound}
\end{equation}

Let
\begin{equation}
 \cZ_1=\sum_{m,n}|x_my_n\mathfrak m(m,n)|.
 \label{eq:intro-Z1}
\end{equation}
The full physical zero mode is therefore bounded by
\begin{equation}
 J\cZ_1\left\{(\log X)^{-A}
              +X^{-s+\kappa+o(1)}\right\}.
 \label{eq:intro-zero-mass-bound}
\end{equation}
The factor \(\cZ_1\) is deliberate.  From the soft estimate
\(\cZ_1\ll Q^2X^\epsilon\), one may close the large-\(g\) term by
choosing \(\epsilon<s-\kappa\), but one may not convert
\(X^\epsilon(\log X)^{-A}\) into logarithmic saving.  Under the
stronger and independently checkable projective hypothesis
\begin{equation}
 \cZ_1\ll Q^2(\log X)^K
 \label{eq:intro-polylog-mass}
\end{equation}
for a fixed \(K\), the exponent \(A\) can be chosen after \(B,K\), and
\eqref{eq:intro-zero-mass-bound} is
\(O_B(XQ(\log X)^{-B})\).  This distinction is part of the theorem,
not an \(X^\epsilon\)-notation convention.  If
\(\widehat F(0)=0\), the entire mass-sensitive term vanishes and the
soft row estimates suffice.

\subsection{A genuine passage beyond
\texorpdfstring{\(J\)}{J}}

Specialize to
\begin{equation}
 S=R=X^{1/2-\delta+o(1)},\qquad T=RX^\xi.
 \label{eq:intro-Rxi}
\end{equation}
Then \eqref{eq:intro-asym-saving} becomes
\begin{equation}
 \eta_{\rm asym}(\beta,\delta,\xi)=
 \min\left\{
 1-\frac{3\beta}{2},
 \frac{\beta+4\delta-1}{2}-\xi,
 \frac{6\delta-\beta-5\xi}{4}
 \right\}.
 \label{eq:intro-R-saving}
\end{equation}
The old Poisson face is
\[
 \xi_{\rm P}=\frac12+\delta-\beta,
 \qquad T=J\ \Longleftrightarrow\ \xi=\xi_{\rm P}.
\]
Suppose
\begin{equation}
 \beta<\frac23,\qquad \beta<\frac12+\delta,
 \label{eq:intro-beta-conditions}
\end{equation}
with fixed margins.  Then the strict interval
\begin{equation}
 \boxed{
 \frac12+\delta-\beta<\xi<
 \min\left\{
 \frac{\beta+4\delta-1}{2},
 \frac{6\delta-\beta}{5}
 \right\}}
 \label{eq:intro-passage-window}
\end{equation}
is a certified one-sided passage beyond \(J\).  A convenient nonempty
region is
\begin{equation}
 \delta>\frac1{10},\qquad
 \frac58-\frac\delta4<\beta<
 \min\left\{\frac23,\frac12+\delta\right\}.
 \label{eq:intro-nonempty-region}
\end{equation}
For example, \((\delta,\beta,\xi)=(0.15,0.62,0.04)\) gives
\[
 R=X^{0.35},\qquad J=X^{0.38},\qquad T=X^{0.39},
 \qquad \eta_{\rm asym}=0.02.
\]

\subsection{Scope and organization}

The theorem treats the incremental mixed shell
\([1,S]\times(S,T]\), and by transposition its mirror.  Under
\eqref{eq:intro-polylog-mass} it closes that calibrated increment,
including physical frequency zero.  Under only the generic soft mass
it proves the fixed-power nonzero theorem and the exact
mass-sensitive zero estimate \eqref{eq:intro-zero-mass-bound}; this is
not silently promoted to full closure.  The result does not re-prove
the old base square, control the both-ultra region, restore every
endpoint of the TPC-18 residual \citep{WangTPC18}, or produce a
positive prime-pair main term.

\Cref{sec:interface-prefixes} fixes the exact prefix and row
interfaces.  \Cref{sec:asymmetric-four-channel} explains the
four-channel cancellation behind \eqref{eq:intro-joint-minus-drift}.
\Cref{sec:zero-shell} proves \eqref{eq:intro-zero-bound}.
\Cref{sec:asymmetric-conductor} extends the conductor tensor to a
rectangle, and \cref{sec:exponent-passage} proves the ledger and
passage window.  Exact splicing, the failed pointwise shortcut, and
the scope ledger are in \cref{sec:splicing-scope}.  The certificate and
conclusion are in \cref{sec:certificate-conclusion}.
